\documentclass[11pt]{article}
\usepackage[a4paper,margin=1in]{geometry}
\usepackage[T1]{fontenc}
\usepackage{lmodern,microtype}
\usepackage{amsmath,amssymb,amsthm,mathrsfs,booktabs,array,graphicx}
\usepackage[colorlinks=true,allcolors=blue]{hyperref}
\newtheorem{theorem}{Theorem}[section]
\newtheorem{proposition}[theorem]{Proposition}
\newtheorem{corollary}[theorem]{Corollary}
\newtheorem{lemma}[theorem]{Lemma}
\theoremstyle{definition}\newtheorem{definition}[theorem]{Definition}
\theoremstyle{remark}\newtheorem{remark}[theorem]{Remark}
\DeclareMathOperator{\ord}{ord}
\DeclareMathOperator{\Div}{Div}
\DeclareMathOperator{\Norm}{N}
\title{Tagged Prime-Ideal Packets and the Rational Pushforward\\
Obstruction for H\'enon Monodromy}
\author{Liang Wang$^{*1}$\\
\small $^{1}$School of Artificial Intelligence and Automation, Huazhong University of Science and Technology\\
\small Wuhan 430074, P.R. China\\
\small $^*$Corresponding author}
\date{Preprint, August 2026}
\hypersetup{pdfauthor={Liang Wang},pdftitle={Tagged Prime-Ideal Packets and the Rational Pushforward Obstruction for Henon Monodromy}}

\begin{document}
\maketitle

\begin{abstract}
Full multiplier-field norms of primitive cyclotomic H\'enon packets are
squares, but their inversion-fixed trace-field factors and prime ideals
remain source-native arithmetic objects.  We construct an exact
finite-cutoff assembly of those survivors.  Each atom retains the signed
primitive orbit, cyclotomic index, trace-field prime ideal, residue degree,
and ideal valuation.  The rational norm is an exact pushforward of this
tagged divisor.  Moreover, at residue characteristic not dividing the
cyclotomic index, every supporting multiplier-field prime sees the unstable
multiplier with exact residue order equal to that index.  An exact
certificate for signed primitive H6 periods one, three, and four and indices
$3\le n\le20$ contains 125 tagged atoms over 95 rational primes.  The free
rational pushforward therefore has kernel rank 30.  Explicit collisions
show that the same rational prime carries different orbit labels, different
indices, different trace primes, and even incompatible exact residue orders;
$p=29$, for example, carries orders 7, 14, and 15.  Thus a rational-prime-
only ledger is not a lossless H6 Euler clock.  The obstruction is scoped:
weighted scalar statistics may remain useful, but an all-orbit trace first
requires a convergent vector-valued tagged assembly.  No analytic
continuation or Hilbert--P\'olya operator is claimed.
\end{abstract}

\section{Introduction}

Periodic-orbit constructions often turn local determinant data into scalar
Euler factors.  For the certified area-preserving H6 H\'enon survivor, this
strategy encounters a sequence of algebraic obstructions.  The primitive
return matrices have reciprocal algebraic-unit multipliers, but no common
real power maps the exact primitive controls to rational primes.  Passing to
cyclic resultants does produce integral sequences; taking the full
multiplier-field norm, however, counts the inversion pair twice and yields a
square.  The smallest surviving arithmetic objects are therefore the
unit-normalized cyclotomic values in the fixed trace field and their
prime-ideal factorizations.

Cyclic resultants and their recurrences are well developed for a fixed
polynomial \cite{hillar2005cyclic,hillarlevine2007recurrences}.  Primitive
divisor theorems also give strong information for Lucas, Lehmer, and certain
quadratic Lehmer--Pierce sequences \cite{bhv2001,flatters2009}.  Arithmetic
dynamics supplies related Zsigmondy phenomena for iterates of a fixed
rational map \cite{ingramsilverman2009,silverman2012}.  None of these results
provides a common divisor ledger for varying signed H6 monodromy fields and
pressure-weighted primitive orbits.

This paper takes one deliberately finite but exact step.  We define a free
tagged divisor group whose basis remembers an orbit, a cyclotomic index, and
a trace-field prime ideal.  We prove that rational norms are exact
residue-degree pushforwards of this object.  We then attach an intrinsic
clock: away from residue characteristics dividing the index, every
multiplier-field prime above a packet atom gives exact multiplier order.
Finally, an executable H6 certificate proves that forgetting to rational
primes is noninjective.

The result has two complementary messages.  Positively, the phrase
``retain the prime ideals'' becomes a concrete finite data type with exact
functoriality.  Negatively, a bare rational prime cannot be a lossless local
label: it can support incompatible source clocks.  This does not rule out
derived scalar statistics.  It says that scalar pushforward must be treated
as a lossy map, not as an identification.

\section{Signed H6 fields and half packets}

Consider
\[
H_6(q,p)=(1-6q^2-p,q),
\qquad
J(q)=\begin{pmatrix}-12q&-1\\1&0\end{pmatrix}.
\]
Every chronological return matrix has determinant one.  For a primitive
orbit $\gamma$, let $\lambda_\gamma$ denote the actual signed unstable
eigenvalue.  Its reciprocal is the second eigenvalue.  The sign is part of
the source data: the exact period-three branch is negative.

The three controls used below have multiplier polynomials
\begin{align}
f_1(X)&=X^4-4X^3-22X^2-4X+1,\label{eq:f1}\\
f_3^-(X)&=X^4+76X^3-7374X^2+76X+1,\label{eq:f3}\\
f_4(X)&=X^2-578X+1.\label{eq:f4}
\end{align}
They are reciprocal and monic with constant term one.  Write
$K_\gamma=\mathbb Q(\lambda_\gamma)$, let $\iota$ be inversion, and set
$F_\gamma=K_\gamma^{\langle\iota\rangle}$.  Their fixed trace fields and
trace values $t_\gamma=\lambda_\gamma+\lambda_\gamma^{-1}$ are
\[
\begin{array}{c|c|c}
\text{primitive period}&F_\gamma&t_\gamma\\\hline
1&\mathbb Q(\sqrt7)&2+2\sqrt7\\
3&\mathbb Q(\sqrt5)&-38-42\sqrt5\\
4&\mathbb Q&578.
\end{array}
\]

For $n>2$, define the inversion-fixed half packet
\begin{equation}\label{eq:beta}
\beta_{\gamma,n}
=\lambda_\gamma^{-\varphi(n)/2}\Phi_n(\lambda_\gamma).
\end{equation}
Since $\varphi(n)$ is even, cyclotomic reciprocity shows that
$\iota(\beta_{\gamma,n})=\beta_{\gamma,n}$.  The multiplier is a unit, so
$\beta_{\gamma,n}\in\mathcal O_{F_\gamma}$.  This is the trace-field
survivor of the full-field square norm.

\section{Finite tagged assembly}

Let $S$ be a finite set of pairs $(\gamma,n)$ with $n>2$.  For every
supporting trace-field prime define a formal source atom.

\begin{definition}[Tagged packet divisor]
Let
\[
\mathcal A_S=\{(\gamma,n,\mathfrak q):
(\gamma,n)\in S,\ \mathfrak q\mid(\beta_{\gamma,n})\}
\]
and
\[
\Div_{\mathrm{tag}}(S)
=\bigoplus_{(\gamma,n,\mathfrak q)\in\mathcal A_S}
\mathbb Z[\gamma,n,\mathfrak q].
\]
The finite tagged divisor is
\begin{equation}\label{eq:tagdiv}
\mathscr D_S=\sum_{(\gamma,n)\in S}
\sum_{\mathfrak q}v_{\mathfrak q}(\beta_{\gamma,n})
[\gamma,n,\mathfrak q].
\end{equation}
\end{definition}

Here canonical means canonical relative to the fixed signed monodromy and
the finite source set $S$.  We do not claim universal minimality among all
possible encodings.

Let $p=\mathfrak q\cap\mathbb Z$ and let
$f(\mathfrak q/p)$ be the residue degree.  Define
\begin{equation}\label{eq:push}
\nu_S:[\gamma,n,\mathfrak q]\longmapsto
f(\mathfrak q/p)[p]
\end{equation}
from the tagged group to the free rational-prime divisor group.

\begin{theorem}[Exact finite assembly and norm pushforward]
\label{thm:assembly}
The divisor \eqref{eq:tagdiv} is a lossless finite record of the source
prime-ideal factorizations.  For every $(\gamma,n)\in S$,
\begin{equation}\label{eq:normpush}
\operatorname{div}_{\mathbb Z}
\left|\Norm_{F_\gamma/\mathbb Q}(\beta_{\gamma,n})\right|
=\sum_{\mathfrak q\mid(\beta_{\gamma,n})}
v_{\mathfrak q}(\beta_{\gamma,n})
f(\mathfrak q/p)[p].
\end{equation}
Consequently the rational divisor of the product of all packet norms is
$\nu_S(\mathscr D_S)$.
\end{theorem}
\begin{proof}
Unique factorization of nonzero ideals gives
$(\beta_{\gamma,n})=\prod_{\mathfrak q}\mathfrak q^{v_{\mathfrak q}}$.
The ideal norm of $\mathfrak q$ is $p^{f(\mathfrak q/p)}$.  Taking norms and
then rational prime divisors proves \eqref{eq:normpush}.  The source tags are
basis coordinates in a finite free group, so they are retained before the
pushforward.  Linearity proves the final assertion.
\end{proof}

\begin{remark}[Where the content lies]
The retention of the tags is definitional once
$\Div_{\mathrm{tag}}(S)$ is chosen.  The substantive assertions are that the
H6 half packets canonically supply its coefficients, that residue degrees
give the exact norm map \eqref{eq:normpush}, and that the same atoms carry
the good-characteristic order in Theorem~\ref{thm:order}.  No novelty or
minimality claim is made for free divisor groups in general.
\end{remark}

\section{The intrinsic good-characteristic clock}

A trace-field atom supports primes in the multiplier field.  These two
prime levels must be kept distinct.

\begin{theorem}[Exact residue order]
\label{thm:order}
Let $\mathfrak q\mid(\beta_{\gamma,n})$ and let $\mathfrak P$ be any prime
of $K_\gamma$ above $\mathfrak q$, of residue characteristic $p$.  If
$p\nmid n$, then the reduction of $\lambda_\gamma$ in
$(\mathcal O_{K_\gamma}/\mathfrak P)^\times$ has exact multiplicative order
$n$.
\end{theorem}
\begin{proof}
The unit factor in \eqref{eq:beta} implies
$(\beta_{\gamma,n})\mathcal O_{K_\gamma}
=(\Phi_n(\lambda_\gamma))$.  Hence
$\Phi_n(\lambda_\gamma)\in\mathfrak P$.  Write $\bar\lambda$ for the
reduction.  Since $p\nmid n$, $X^n-1$ is separable in characteristic $p$.
In the reduced factorization
\[
X^n-1=\prod_{d\mid n}\Phi_d(X),
\]
the monic factors are therefore pairwise coprime.  Now
$\Phi_n(\bar\lambda)=0$ implies $\bar\lambda^n=1$.  If its order were a
proper divisor $d<n$, then $\bar\lambda$ would also be a root of
$X^d-1=\prod_{e\mid d}\Phi_e(X)$, contradicting that coprimality.  Its
order is exactly $n$.
\end{proof}

\begin{remark}[Bad characteristic]
When $p\mid n$, exact order $n$ is impossible in a finite residue field:
the multiplicative group has order $p^f-1$, which is prime to $p$.  The
reduced cyclotomic factors also need not remain coprime.  Such atoms are
retained in the divisor ledger, but their actual smaller order is not
asserted or computed by Theorem~\ref{thm:order}.
\end{remark}

Thus $n$ is not merely a bookkeeping index at good characteristic.  It is an
intrinsic finite-field clock.  Any pushforward that merges atoms with
different $n$ can merge different exact clocks.

\section{The rational pushforward obstruction}

The norm map in Theorem~\ref{thm:assembly} is exact but need not be
injective.  The following rank formula makes the information loss explicit.

\begin{proposition}[Free kernel rank]
\label{prop:kernel}
Let $m=\#\mathcal A_S$, and let $r$ be the number of distinct rational
primes below atoms in $\mathcal A_S$.  Then
\[
\operatorname{rank}_{\mathbb Z}\ker\nu_S=m-r.
\]
In particular, $m>r$ implies noninjectivity.
\end{proposition}
\begin{proof}
For a fixed rational prime $p$, suppose $m_p$ tagged atoms lie above it.
The $p$-block of \eqref{eq:push} is
\[
\mathbb Z^{m_p}\longrightarrow\mathbb Z,
\qquad (x_i)\longmapsto\sum_i f_i x_i,
\]
with all $f_i>0$.  Its rank is one and its free kernel has rank $m_p-1$.
Summing these independent blocks over the $r$ rational primes gives
$\sum_p(m_p-1)=m-r$.
\end{proof}

Noninjectivity is a scoped obstruction.  It does not prohibit a useful
weighted pushforward; it prohibits calling the untagged rational divisor a
lossless reconstruction of the source packet.

\section{Exact H6 certificate}

The certificate computes the half packets for the signed controls
\eqref{eq:f1}--\eqref{eq:f4} and indices $3\le n\le20$.  It reconstructs the
trace polynomials, uses exact integral bases, factors ramified, inert, and
split trace primes, and checks the multiplier order in every supporting
finite-field factor.  No floating-point approximation or external prime
table is used.

The finite ledger contains
\[
54\ \text{packet rows},\qquad
125\ \text{tagged atoms},\qquad
95\ \text{distinct rational primes}.
\]
Proposition~\ref{prop:kernel} therefore gives
\[
\operatorname{rank}\ker\nu_S=125-95=30.
\]
This is not merely a signed-group artifact.  At $p=109,n=11$, three distinct
effective basis atoms all have residue degree one and hence each maps to the
same effective divisor $[109]$.  Noninjectivity is already visible on the
positive atom set.
There are 105 good-characteristic atoms with exact order and 20
bad-characteristic atoms deliberately left uncertified.  All 54 rational
norm factorizations reconstruct exactly, including 30 inherited half-norm
crosschecks.  The canonical core digest is
\begin{center}
\small\texttt{9cc40f4b8ed5f54e36c70cb1db15b509837aee264e86bcf39395f827e59153ba}.
\end{center}

Table~\ref{tab:collisions} shows decisive source collisions.

\begin{table}[t]
\centering
\caption{Selected tagged atoms.  The order column is the exact multiplier
residue order; all displayed rows are good characteristic.}
\label{tab:collisions}
\small
\resizebox{\textwidth}{!}{\begin{tabular}{rrrrrrr}
\toprule
$p$ & orbit & $n$ & trace prime ideal & $f$ & $v_{\mathfrak q}$ & order \\
\midrule
29 & $\gamma_1$ & 15 & (29,w-6) & 1 & 1 & 15 \\
29 & $\gamma_3^{-}$ & 7 & (29,w-24) & 1 & 1 & 7 \\
29 & $\gamma_3^{-}$ & 14 & (29,w-6) & 1 & 1 & 14 \\
109 & $\gamma_1$ & 11 & (109,w-15) & 1 & 1 & 11 \\
109 & $\gamma_1$ & 11 & (109,w-94) & 1 & 1 & 11 \\
109 & $\gamma_3^{-}$ & 11 & (109,w-99) & 1 & 1 & 11 \\
131 & $\gamma_1$ & 12 & (131,w-111) & 1 & 1 & 12 \\
131 & $\gamma_3^{-}$ & 5 & (131,w-120) & 1 & 1 & 5 \\
38039 & $\gamma_3^{-}$ & 19 & (38039,w-22136) & 1 & 1 & 19 \\
38039 & $\gamma_4$ & 13 & (38039) in Z & 1 & 1 & 13 \\
\bottomrule
\end{tabular}
}
\end{table}

The table exposes three independent tag requirements on this certificate.
At $p=109,n=11$, periods one and three coexist, so the orbit tag is needed.
Within period one, two distinct split trace primes coexist, so the
prime-ideal tag is needed.  At $p=29$, the certified orders 7, 14, and 15
coexist, so retaining only the rational prime loses the cyclotomic index and
clock.  Further multi-order collisions occur at
\[
p\in\{11,19,29,79,131,307,38039\}.
\]
The signed period-three branch is separately mutation-tested: replacing it
by the positive modulus changes its packet.

These controls do not prove that our key is categorically minimal.  They
prove that the displayed orbit, index, prime-ideal, and sign coarsenings lose
actual information in the exact H6 data.

The certificate is finite and is not evidence for an asymptotic density of
collisions or of primitive rational divisors.  Its role is sharper and
narrower: it is an exact counterexample to injectivity of the proposed
rational-prime-only source map.

\section{Source boundary and the next theorem}

Hillar's cyclic-resultant framework and the Hillar--Levine recurrence theory
concern one fixed polynomial \cite{hillar2005cyclic,hillarlevine2007recurrences}.
Flatters' theorem supplies eventual primitive rational divisors for the
period-four real-quadratic Lehmer--Pierce sequence \cite{flatters2009}, while
the Bilu--Hanrot--Voutier theorem treats Lucas and Lehmer numbers
\cite{bhv2001}.  Ingram--Silverman and related dynamical-Zsigmondy work treat
orbits of a fixed rational map \cite{ingramsilverman2009,silverman2012}.
These are strong positive analogues, but none accepts the present family of
varying signed H6 monodromy fields and pressure weights.

The next theorem must therefore begin with the tagged source object.  A
schematic target is a pressure-filtered divisor measure
\[
\mathfrak D_T
=\sum_{\substack{\gamma\ \mathrm{primitive}\\\ell_P(\gamma)\le T}}
W_T(\gamma)
\sum_{n\ge3}\sum_{\mathfrak q\mid(\beta_{\gamma,n})}
v_{\mathfrak q}(\beta_{\gamma,n})[\gamma,n,\mathfrak q].
\]
One must specify a common vector-valued topology, derive $W_T$ from the H6
pressure clock, and prove convergence before applying a rational norm map.
No such theorem is supplied here.

The strict Route-A tuple is
\[
(A1_{\rm WEAK},A2_{\rm FAIL},A3_{\rm FAIL},A4_{\rm FORMAL\_HINT}),
\qquad \texttt{ROUTE\_A\_EXPLORATORY}.
\]
There is no global Fredholm determinant, continuation theorem, functional
equation, Hilbert space, or self-adjoint operator.  Route B is not
authorized.

\paragraph{Reproducibility.}
The certificate is regenerated by
\texttt{bash code/run\_c50.sh}.  The script uses exact integer, ideal-basis,
and finite-field polynomial arithmetic, runs eleven unit and adversarial
tests, and emits both the canonical JSON ledger and
Table~\ref{tab:collisions}.  It reads no prime table or Riemann-zero table.

\section{Conclusion}

Prime-ideal packets can be assembled across a finite signed H6 source set
without losing their provenance.  Rational norms are exact pushforwards, and
good-characteristic packet atoms carry the cyclotomic index as an exact
residue-order clock.  The same exactness reveals the obstruction: the
rational-prime-only map has a positive free kernel and merges incompatible
clocks.  The mathematically honest next bridge is therefore vector-valued
and pressure weighted.  Scalarization, if useful, must come after that
assembly theorem rather than replace it.

\bibliographystyle{plain}
\bibliography{references}
\end{document}
